%-*-tex-*-
\ifundefined{writestatus} \input status \relax \fi %
\chcode{over}
\def\cqu{\cquote{There is no one on earth who does what is right
all the time and never makes a mistake.}{Ecclesiastes Chap. 7, Verse 20}}

\chapterhead{over}{OVERFULL BOXES\cr and\cr OTHER DISASTERS}
Sometimes \tex\ and \intex\ have trouble making a {\it best} decision that is
good enough. The result of this is a message that there is an {\it overfull}
or {\it underfull} box. Chapter \ref{parcont} suggests increasing the
value of |\tolerance| and |\pretolerance| so that \tex\ will accept worse
looking lines. This is fine except that the result is less than beautiful. 

\tex\ indicates an overfull box by putting a big ugly black  line at the end.
This book has a few overfull boxes to demonstrate how they look. This chapter
suggests ways of both avoiding and hiding overfull boxes.

\shead{overcomform}{Command Forms}
\bshortcomlist
\pri\@|\hfuzz [=] <dimen>|&\tex\ tells of overfull horizontal boxes only if
they are overfull by more than |<dimen>|. {\it Plain} default is 0.1pt. \cr
\pri\@|\overfullrule [=] <dimen>|&\tex\ places a rule or line of width
|<dimen>| at the end of an overfull box and of the same height as the box.\cr
\pri\@|\vfuzz [=] <dimen>|&\tex\ tells of overfull vertical boxes only if
they are overfull by more than |<dimen>|. {\it Plain} default is 0.1pt. \cr
\eshortcomlist

\shead{action}{Overfull Boxes --- Action}
The recommended way to handle an overfull box is to rewrite the passage that
is causing the offence. This will generally work quite well if the box is an
isolated event. It has the feature of not destroying the beauty of the
document.

However, if you want to hear about overfull boxes but don't want the black
marks let |\overfullrule = 0pt|. If you don't want to hear about them,
increase the size of |\hfuzz| and |\vfuzz|.
|\finalversion| increases |\hfuzz| and obliterates the black marks. 

Overfull |vboxes| are quite rare, as are underfull |hboxes|. If there is an
underfull |hbox|, check it very carefully. The most likely error is not
putting in glue to fill in the box. An example would be using |\line|
instead of |\centerline|. 

Underfull |vboxes| are not as rare. They result when the |vboxes| 
on the page are quite large and the glue between them is not very stretchy.
Pages with figures, tables or multicolumn format are the most likely
candidates. The only solution may be a rearrangement of material in the
document, forced filling of the bottom of the page. \tex\ provides a
\@|\raggedbottom| command that puts glue at the bottom of a short page to
prevent it from being stretched out. Whether this is an adequate solution to
the problem is a matter of aesthetic debate.
In general, information about over and under full boxes is desirable. 

\begingroup \bf
When \tex\ indicates that there is an underfull box, it stretches, unless
|\raggedbottom| is active,  all stretchable glue. This may result in pages
that look strange. Only glue between paragraphs, before and after display
mathematics, between all list items, and before section \dots\ heads is
stretchable. Interline glue is not stretchable either in {\it Plain\/} or
\intex. In extreme cases (badness greater 10000) the page will have very
disturbing spacing indeed. 
\par\endgroup


\ejectpage
\done